% =============================================================================
% Kapitel 2: Theoretische Grundlagen
% =============================================================================
\chapter{Theoretische Grundlagen}
\label{ch:grundlagen}

\todo{Einleitenden Text für das Kapitel verfassen. Was wird in diesem Kapitel behandelt und warum ist es relevant?}

In diesem Kapitel werden die theoretischen Grundlagen für die vorliegende Arbeit festgehalten, um ein flächendeckendes Verständnis für die folgenden Kapitel zu schaffen. Es werden relevante Themenbereiche erläutert, der aktuelle Stand der Forschung zusammengefasst und verwandte Arbeiten vorgestellt. Ziel ist es, die eigene Arbeit von bereits existierenden Ansätzen und Erkenntnissen abzugrenzen und die Relevanz der eigenen Forschung zu unterstreichen.

% -----------------------------------------------------------------------------
\section{Erfassen von Mobilitätsdaten}
\label{sec:themenbereich1}



% -----------------------------------------------------------------------------
\section{Verarbeitung von Mobilitätsdaten}
\label{sec:themenbereich2}

% -----------------------------------------------------------------------------
\section{Bewertung der Datenqualität}
\label{sec:themenbereich3}



% -----------------------------------------------------------------------------
\section{Stand der Forschung}
\label{sec:stand_forschung}

Die Forschung zur Prognose individueller Fahrzeugnutzung lässt sich grob in
mehrere Stränge gliedern, die sich in Zielgröße, Methodik und Datengrundlage
unterscheiden. Tabelle~\ref{tab:forschungsstraenge} fasst diese Stränge als
Übersicht zusammen; die folgenden Unterabschnitte vertiefen sie jeweils und
arbeiten am Ende die für diese Arbeit zentrale Forschungslücke heraus.

\begin{table}[htbp]
  \centering
  \begin{tabularx}{\textwidth}{@{}>{\raggedright\arraybackslash}p{3.2cm} X X X@{}}
    \toprule
    \textbf{Forschungsstrang} & \textbf{Typische Methoden} &
    \textbf{Datengrundlage} & \textbf{Offene Punkte} \\
    \midrule
    Mobilitäts- und Nutzungsforecast & Statistische Modelle (ARIMA),
    Random Forest, neuronale Netze (LSTM) & Reale Telematik- und
    Bewegungsdaten & Generalisierung über Fahrzeuge und Quellen hinweg \\
    \addlinespace
    Driving-/Parking-Klassifikation & Random Forest, SVM, Gradient
    Boosting & EV-Flottendaten, stündliche Profile & Einfluss der
    Datenqualität auf die Klassifikationsgüte \\
    \addlinespace
    EV-Last- und Ladeprognose & Gradient Boosting, Deep Learning &
    Synthetische Daten (z.\,B.\ emobpy), reale EV-Daten & Vergleichbarkeit
    synthetischer und realer Datengrundlagen \\
    \addlinespace
    Unsicherheits- und Quantilsprognose & Quantilsregression,
    Quantile-Random-Forest & Heterogene Mobilitätsdaten & Kalibrierung
    individueller Vorhersageintervalle (P10/P90) \\
    \bottomrule
  \end{tabularx}
    \caption{Zentrale Forschungsstränge im Umfeld der individuellen
    Fahrzeugnutzungsprognose.}
    \label{tab:forschungsstraenge}
\end{table}
\todo{Aktuellen Stand der Forschung zusammenfassen. Welche verwandten Arbeiten gibt es? Was wurde bereits untersucht? Pro Tabellenzeile eine Subsection mit synthetisierter Quellenlage anlegen, abschließend Forschungslücke formulieren.}

\subsection{Mobilitäts- und Nutzungsforecast}
\label{subsec:mobilitaetsforecast}

Ein zentraler Befund dieses Strangs betrifft die Heterogenität individuellen
Fahrverhaltens. \textcite{ji2023rethinking} werten über vier Millionen Fahrten
von 3743 privaten Fahrzeugführern in sechs US-Metropolregionen aus und zeigen,
dass sich die Fahrer in fünf deutlich verschiedene Typen gliedern lassen
(\Cref{tab:fahrertypen}), die sich vor allem in ihrer Pendelregelmäßigkeit und
im Anteil nicht-arbeitsbezogener Fahrten unterscheiden. Diese Heterogenität
erklärt, warum ein einzelnes, übergreifendes Modell dem Individualverkehr nur
begrenzt gerecht wird, und motiviert eine individuelle bzw. quellenspezifische
Modellierung.

\begin{table}[htbp]
  \centering
  \begin{tabularx}{\textwidth}{lX}
    \toprule
    \textbf{Fahrertyp} & \textbf{Charakterisierung} \\
    \midrule
    Homebodies & Wenigfahrer, die nur kurz und selten das Haus verlassen; geringe Mobilität. \\
    Gig Drivers & Fahrer mit hohem, auftragsgetriebenem Fahraufkommen (z.\,B.\ Liefer- und Fahrdienste) ohne festen Arbeitsort. \\
    Movers & Individuen mit breit gestreutem Fahrverhalten und vielen wechselnden Zielen, ohne klares Muster. \\
    Typical Drivers & Durchschnittsfahrer mit einer Mischung aus Arbeits- und Freizeitfahrten und mittlerer Regelmäßigkeit. \\
    Work-focused Commuters & Berufspendler mit stark regelmäßigem Pendelmuster, wenig nicht-arbeitsbezogenen Fahrten und hoher Vorhersagbarkeit. \\
    \bottomrule
  \end{tabularx}
  \caption{Fünf Fahrertypen des Individualverkehrs nach \textcite{ji2023rethinking}.}
  \label{tab:fahrertypen}
\end{table}

% -----------------------------------------------------------------------------
\section{Verwandte Arbeiten}
\label{sec:verwandte_arbeiten}

\todo{Verwandte Arbeiten vorstellen und voneinander abgrenzen. Wie unterscheidet sich die eigene Arbeit?}



\textcolor{red}{
    \begin{itemize}
        \item Forschungsarbeit/SUMO
        \item ähnliche Studien
        \item 
    \end{itemize}
}

% -----------------------------------------------------------------------------
\section{Zusammenfassung}
\label{sec:grundlagen_zusammenfassung}

\todo{Kurze Zusammenfassung der wichtigsten Erkenntnisse aus diesem Kapitel.}
